\section{Phụ lục}

\subsection{Tóm tắt ảnh đã tạo bằng code}
\begin{itemize}[leftmargin=2em]
    \item \texttt{img/generated/overview\_top\_countries.png}: top quốc gia theo số lượng phim.
    \item \texttt{img/generated/overview\_top\_languages.png}: top ngôn ngữ theo số lượng phim.
    \item \texttt{img/generated/quality\_rating\_distribution.png}: phân bố phim theo nhóm điểm đánh giá.
    \item \texttt{img/generated/finance\_top\_roi.png}: top 10 phim có ROI cao nhất.
    \item \texttt{img/generated/creative\_genre\_mix.png}: tác động của số lượng thể loại đến độ phổ biến và điểm trung bình.
\end{itemize}

\subsection{Ảnh cần bổ sung thủ công}
\begin{enumerate}[leftmargin=2em]
    \item Ảnh chụp toàn bộ dashboard ở trạng thái mặc định.
    \item Ảnh tab Toàn cảnh khi hiển thị bản đồ, treemap và line chart theo năm.
    \item Ảnh tab Chất lượng khi hiển thị donut, bar chart và heatmap.
    \item Ảnh tab Tài chính khi hiển thị top ROI, line chart và bubble chart.
    \item Ảnh tab Sáng tạo khi hiển thị scatter đạo diễn, bar diễn viên và biểu đồ số thể loại.
\end{enumerate}

\subsection{Gợi ý thao tác lấy ảnh}
Chạy lệnh \texttt{streamlit run src/main.py}, giữ bộ lọc mặc định để ảnh trong báo cáo thống nhất với phần mô tả. Nếu cần ảnh sắc nét cho báo cáo PDF, nên chụp ở độ phân giải cao và lưu trong thư mục \texttt{img/} rồi thay placeholder tương ứng trong file \texttt{content/how-dashboard.tex}.
